\section{The raw counit for categories already of connective-chain form}

\subsection{F-SM39: coinduced disks remove the scalar-dimension restriction}

Let \(k\) be an arbitrary commutative ring, and let \(R\) be a
small \(k\)-linear category such that every \(R(a,b)\) is a
projective \(k\)-module; finite generation is not assumed.
Put \(\mathcal N=\operatorname{Mod}\!-\!R\), with right-module
convention, and suppose that
\[
 \varphi_{\mathcal N}\dashv\kappa_{\mathcal N}\dashv
 \psi_{\mathcal N},\qquad
 \kappa_{\mathcal N}:\operatorname{Mod}(k)\longrightarrow\mathcal N
                                                               \tag{SM39.1}
\]
are \(k\)-linear adjunctions with \(\kappa_{\mathcal N}\) exact.
Assume the canonical derived counit at the scalar ring is a
quasi-isomorphism:
\[
       L\varphi_{\mathcal N}(\kappa_{\mathcal N}k)
                         \longrightarrow k[0].              \tag{SM39.2}
\]
Set
\[
 \begin{aligned}
 \mathcal M&=\operatorname{Ch}_{\ge0}(\mathcal N),\\
 \kappa_{\mathcal M}(U)&=\kappa_{\mathcal N}(U)[0],\\
 \varphi_{\mathcal M}(X)&=\varphi_{\mathcal N}(H_0X).
 \end{aligned}                                               \tag{SM39.3}
\]
All complexes in this statement are homological.

Use the standard compact projective generators of \(\mathcal M\):
if \(P_a=R(-,a)\), take
\[
 Q_{a,0}=P_a[0],\qquad Q_{a,m}=D_m(P_a)\quad(m\ge1),
                                                               \tag{SM39.4}
\]
where \(D_m(A)\) has \(A\) in degrees \(m,m-1\) and identity
differential.  Let \(S\) be their full small \(k\)-linear
category, identify \(\mathcal M\) with right \(S\)-modules,
and let \(D(s)=\varphi_{\mathcal M}(Q_s)\) be the left
coefficient module.  Choose any resolution
\(B_\bullet\to D\) by projective left \(S\)-modules.
Form the exact adjunction of BR-13405:
\[
 \begin{aligned}
 \Lambda(X)&=B_\bullet\otimes_S X,\\
 (GC)(s)&=\operatorname{Hom}_{\operatorname{Ch}_{\ge0}(k)}
                                  (B_\bullet(s),C).
 \end{aligned}                                               \tag{SM39.5}
\]
Then the \emph{ordinary, unreplaced} counit
\[
                    \epsilon_C:\Lambda G(C)\longrightarrow C
                                                               \tag{SM39.6}
\]
is a quasi-isomorphism for every connective complex \(C\).
Consequently this very adjunction induces
\[
 \mathcal M[W^{-1}]\simeq D_{\ge0}(k),\qquad
          W=\Lambda^{-1}(\mathrm{qis}).                     \tag{SM39.7}
\]
No finite injective or pure-injective dimension for scalar
projectives, no finitely generated scalar morphism modules,
and no condition \(D\in\mathrm{FP}_\infty\) is imposed.

\begin{proof}
\emph{1. Generators, exactness, and the derived functor.}
The generators (SM39.4) satisfy
\[
       \operatorname{Hom}_{\mathcal M}(Q_{a,m},X)=X_m(a).
                                                               \tag{SM39.8}
\]
These evaluation functors are exact and preserve coproducts,
so the \(Q_{a,m}\) are compact projective generators.
Presentations by their coproducts give the stated restricted
Yoneda equivalence with \(\operatorname{Mod}\!-\!S\).
Formula (SM39.8) also shows that every \(S(s,t)\) is either
zero or a scalar morphism module of \(R\); in particular it
is projective over \(k\).

The scalar adjunctions of \(\mathcal M\) are
\[
 \varphi_{\mathcal M}\dashv\kappa_{\mathcal M}\dashv
        (\psi_{\mathcal N}\circ\operatorname{ev}_0).
                                                               \tag{SM39.9}
\]
For the first adjunction, a chain map into an object in degree
zero factors uniquely through \(H_0\); for the second, a
chain map from an object in degree zero is an arbitrary map
to the degree-zero term.  The middle functor is exact.
It follows that \(\varphi_{\mathcal M}\) sends projectives
to projective scalar modules.  Thus each \(D(s)\) is
projective over \(k\).

The value \(B_\bullet(s)\to D(s)\) is a projective
scalar resolution of a projective scalar module.  Splitting
the augmentation and then its successive kernels, starting
in degree zero, decomposes \(B_\bullet(s)\) into its
projective degree-zero homology and projective disks.
It is therefore a projective object of
\(\operatorname{Ch}_{\ge0}(k)\).  This proves that \(G\)
is exact, by (SM39.5).  Each \(B_n\) is flat over \(S\),
so \(\Lambda\) is exact as a functor to the abelian category
of complexes.  The right exact coproduct-preserving functor
\(\varphi_{\mathcal M}\) is \(D\otimes_S-\): the two
functors agree on the compact projective generators and
their coproducts, and hence on their cokernel presentations.
Consequently
\[
                 \Lambda(X)\simeq L\varphi_{\mathcal M}(X)
                    \quad\hbox{in }D_{\ge0}(k).              \tag{SM39.10}
\]

We will need the further identification
\[
       L\varphi_{\mathcal M}(X)
                         \simeq L\varphi_{\mathcal N}(X),     \tag{SM39.11}
\]
where the right side applies the derived functor to the
\(\mathcal N\)-complex \(X\).  Here is a finite-diagonal
proof.  Resolve \(X\), as an object of \(\mathcal M\), by
projective chain objects \(A_r\), \(r\ge0\).  A projective
chain object is a retract of a coproduct of the generators
(SM39.4); hence its terms and its \(H_0\) are projective
in \(\mathcal N\), and its positive homology is zero.
The bicomplex \((A_r)_m\) has finite total diagonals.
Exactness of the resolution in every fixed degree gives
\(\operatorname{Tot}(A_\bullet)\to X\) a quasi-isomorphism.
The row augmentations give a second quasi-isomorphism
\[
          \operatorname{Tot}(A_\bullet)
                            \longrightarrow H_0(A_\bullet).
\]
This follows by taking row homology, which is concentrated
in degree zero.  Thus \(H_0(A_\bullet)\) is a connective
projective \(\mathcal N\)-complex representing \(X\), and
applying \(\varphi_{\mathcal N}\) to it computes both sides
of (SM39.11).  The augmentations make the identification
compatible with the canonical scalar counits.  In particular,
if \(X\) is contractible as an outer \(\mathcal N\)-complex,
then
\[
                            L\varphi_{\mathcal M}(X)=0.       \tag{SM39.12}
\]

\emph{2. The coinduced objects are outer disks.}
For \(a\in R\) and a scalar module \(U\), let
\[
             I_a(U)(b)=\operatorname{Hom}_k(R(a,b),U).
\]
It is the right adjoint to evaluation at \(a\) in
\(\mathcal N\).  For every \(m\ge0\), the right adjoint
to degree-\(m\) evaluation on connective homological
complexes is \(D_{m+1}\): a map
\(X\to D_{m+1}(A)\) is uniquely determined by its degree-\(m\)
component \(X_m\to A\).  Therefore the right adjoint to
(SM39.8) is
\[
                        D_{m+1}(I_a(U)).                    \tag{SM39.13}
\]
Equivalently, under the right \(S\)-module description,
\[
 \operatorname{Hom}_k(S((a,m),-),U)
                   \simeq D_{m+1}(I_a(U)).                  \tag{SM39.14}
\]
Every object in (SM39.14) is contractible in the outer chain
direction.  Products of such objects are contractible:
take the product of their contracting homotopies degree by
degree.  This uses products of maps and their homotopy
identities, not exactness of products.  A retract of a
contractible complex is contractible as well: if \(rs=1\)
and \(h\) contracts the ambient complex, then \(rhs\)
contracts the retract.

If \(L\) is any projective left \(S\)-module, it is a
retract of a coproduct of left representables.  Hence
\(\operatorname{Hom}_k(L(-),U)\) is a retract of a product
of (SM39.14), so it is outer contractible.  In particular,
for every \(n\ge1\) and every scalar module \(U\),
\[
 \begin{aligned}
 G(D_n(U))
   &=\operatorname{Hom}_{\mathrm{Ch}}(B_\bullet,D_n(U))\\
   &\simeq\operatorname{Hom}_k(B_{n-1}(-),U)
                 \quad\hbox{is outer contractible}.
 \end{aligned}                                               \tag{SM39.15}
\]
The chain-map identity follows because a map into a disk
is determined by its degree-\((n-1)\) component; the
degree-\(n\) component is its composite with the source
differential.  There is no injectivity assumption on \(U\).

\emph{3. The canonical scalar normalization holds on projectives.}
The functors \(\varphi_{\mathcal N}\) and
\(\kappa_{\mathcal N}\) preserve coproducts, each being a
left adjoint in (SM39.1).  Coproducts are exact in the module
category \(\mathcal N\).  Taking a coproduct of projective
resolutions therefore computes \(L\varphi_{\mathcal N}\)
on a coproduct, and shows that this derived functor preserves
coproducts.  The direct sum of the canonical maps (SM39.2)
gives the canonical quasi-isomorphism for every free scalar
module \(k^{(A)}\).  Taking retracts gives, for every
projective scalar module \(Q\),
\[
       L\varphi_{\mathcal M}(\kappa_{\mathcal M}Q)
                \simeq L\varphi_{\mathcal N}(\kappa_{\mathcal N}Q)
                         \longrightarrow Q[0]               \tag{SM39.16}
\]
as a quasi-isomorphism.

On the other hand, chain maps into \(Q[0]\) factor through
the degree-zero homology of \(B_\bullet\).  Thus
\[
          G(Q[0])=\operatorname{Hom}_k(D(-),Q)
                         \simeq\kappa_{\mathcal M}Q.         \tag{SM39.17}
\]
The counit (SM39.6) for \(Q[0]\) is exactly (SM39.16):
it is evaluation after the augmentation \(B_\bullet\to D\).
In particular it is the canonical counit, not merely some
isomorphism between its source and target in the derived
category.

\emph{4. Test every projective target, without commuting \(G\) with sums.}
Let \(P\) be a projective object of
\(\operatorname{Ch}_{\ge0}(k)\).  The projective-chain
decomposition BR-13402 gives projective scalar modules
\(Q_0,Q_1,Q_2,\ldots\) and an isomorphism
\[
          P\simeq Q_0[0]\oplus\prod_{n\ge1}D_n(Q_n).
                                                               \tag{SM39.18}
\]
The displayed product of disks is also their coproduct:
only two disks contribute in any fixed degree.
Because \(G\) is a right adjoint, it preserves this product.
By (SM39.15), its value on the product of disks is an outer
contractible object \(K\in\mathcal M\).  Therefore
\[
                 G(P)\simeq\kappa_{\mathcal M}Q_0\oplus K.
                                                               \tag{SM39.19}
\]
Equations (SM39.10)--(SM39.12) show that \(\Lambda K\) is
acyclic.  The projection in (SM39.18) and its image under
\(\Lambda G\) are consequently quasi-isomorphisms.  The
counit square for that projection commutes, and its bottom
map is the quasi-isomorphism (SM39.16).  Two-out-of-three
proves that \(\epsilon_P\) is a quasi-isomorphism.
Only a finite biproduct was passed through \(\Lambda\) in
(SM39.19); no product-preservation assertion about
\(\Lambda\) or coproduct-preservation assertion about \(G\)
has entered.

\emph{5. Extend the test to all targets and descend the adjunction.}
For completeness, the projective-target criterion F-SM17
applies directly as follows.  Resolve an arbitrary \(C\),
as an object of the abelian category of connective
complexes, by projective chain objects \(P_r\).
Both the identity functor and \(\Lambda G\) are exact.
The augmented first-quadrant bicomplexes therefore give
quasi-isomorphisms
\[
 \operatorname{Tot}(\Lambda G(P_\bullet))\longrightarrow\Lambda G(C),
 \qquad
 \operatorname{Tot}(P_\bullet)\longrightarrow C.
\]
The map between the total complexes is a quasi-isomorphism
because its map on each row is \(\epsilon_{P_r}\), already
proved to be one.  All total diagonals are finite, so no
infinite-totalization issue or continuity assumption is
needed.  Naturality and two-out-of-three give (SM39.6) for
every \(C\).

If \(f\) is a quasi-isomorphism of scalar complexes,
naturality of these counits shows that \(\Lambda G(f)\)
is a quasi-isomorphism.  Thus \(G\) is homotopical for the
specified weak equivalences.  The triangle identity
\(\epsilon_{\Lambda X}\Lambda(\eta_X)=1\) shows that
each unit \(\eta_X:X\to G\Lambda X\) is a weak
equivalence.  Both ordinary functors therefore descend,
and their unit and counit become isomorphisms, proving
(SM39.7).
\end{proof}

\paragraph{Scope and comparison with the earlier sufficient conditions.}
This proves the literal counit for the resolution adjunction
(SM39.5), not only the existence of some inverse to the
localized realization.  It applies to all locally
\(k\)-projective ringoids \(R\) in (SM39.1), including those
with infinitely generated scalar morphism modules, because
the extra connective-chain structure supplies the explicit
coinduced contractions.  It does not establish this
structure for every category in the original corollary.
The general raw-counit question remains outside this class.

The earlier example F-SM28 is a consistency check: its
\(\operatorname{Fun}(\mathbb N,\operatorname{Ch}_{\ge0}(k))\)
can be written in the form (SM39.3).  Its raw counit was
already covered by F-SM24's finitely generated scalar
morphism case, so it is not counted here as a newly solved
example.  The present proof imposes none of that condition,
F-SM27's \(\mathrm{FP}_\infty\) condition, or F-SM24's
scalar pure-injective-dimension bound.
